
Struktur und Eigenschaft

etwas Philosophie bzw. Psychologie lässt sich nicht vermeiden

klassisch werden alle Arten von Werten zusammengewürfelt

BEISPIEL

An dieser Stelle ist die Erkenntnis wichtig, dass Strukturinformationen
und Eigenschaften unterschiedlich zu verarbeiten sind!
Warum ist diese Behauptung so wichtig und richtig?
Ein Beispiel soll helfen, es zu verdeutlichen.

BEISPIEL: GUI
- stellt DAG dar
- leicht per Schleifen zu verarbeiten, von Wurzel aus "in die Tiefe tauchen"
- Voraussetzung: alle Elemente haben Eigenschaften, die uniform
  verarbeitet werden können, z. B. x, y, width, height, foreground, background
- außerdem gibt es Verweise auf enthaltene Kindelemente (Strukturinformation)

Um beide nicht zu vermischen, sollten sie getrennt gehalten werden:

- name
- type
- model
- properties
- reference

Eigenschaften wie Position, Größe, Farbe bieten Meta-Informationen
über enthaltene Elemente. Sie stellen aber selbst keine Elemente
im Sinne einer hierarchischen Struktur (Mikro-/Makrokosmos) dar.

Auf die gleiche Weise könnte eine Eigenschaft "position" existieren,
die für die Sortierung der Elemente eines Behälters sorgt.
